\section{The arithmetic bound at one shift}\label{sec:assembly}

Fix a mesh point $x$, the sets $T_2,T$, the shift $u_0$ and the
nodes of \S\ref{sec:theta-chain}. Average over a nonempty set of words
$W$ with weights $d_u=\operatorname{rk}N_u>0$. Put $y=x^\varepsilon$.
The source-index loss of \eqref{eq:loss-second-error} enters the
excess through its mean. Write
\[
 \bar e_2=\langle e_{2,u}\rangle,\qquad
 V_M=\max_{u\in W}\bigl(\mu_{\max}(M_u)+A\log T_u\bigr)
              -\langle\mu^M_u+A\log T_u\rangle .
\]

\subsection{The loss at a mesh point}

\begin{lemma}\label{lem:assembly-spread}
Suppose the source lattices satisfy the hypotheses of
Theorem~\ref{thm:loss-source}. Take $h_u=h_*-A\log T_u$ and
$|E_u-mAg|\le e_5$ for every word. The source spread $U$ of
\S\ref{sec:period-slopes} satisfies
\[
                     U\le V_M+\bar e_2+2e_5 .
\]
\end{lemma}

\begin{proof}
Put $\xi_u=E_u-mAg$. In \eqref{eq:period-spread} the constants
$h_*$ and $mAg$ cancel. Since $R_u\ge0$,
\[
 U\le V_M+\bar e_2+\langle\xi_u\rangle-\min_u\xi_u
      \le V_M+\bar e_2+2e_5 .\qedhere
\]
\end{proof}

\begin{proposition}\label{prop:assembly-excess}
Suppose the following hold at $T,u_0$ for all words.
\begin{enumerate}[label=\textup{(\roman*)},nosep]
\item The lifts satisfy the hypotheses of
Corollary~\ref{cor:true-integral-chain} and the endpoint inclusions
of Hypothesis~\ref{hyp:chain-endpoints}. Their inert exponents give
$\Gamma,e_1$ of \eqref{eq:loss-net-gain} and \eqref{eq:loss-first-error}.
Their source lattices satisfy the hypotheses of Theorem~\ref{thm:loss-source}.
\item The period maps and their metrics satisfy the hypotheses of
Theorem~\ref{thm:height-account} and Corollary~\ref{cor:bessel-mean},
with $h_u=h_*-A\log T_u$ from the height theorem.
\item The auxiliary primes and base lattices satisfy the hypotheses
of Propositions~\ref{prop:height-prime} and~\ref{prop:height-base}.
The finite exponents of these kernels satisfy the hypotheses of
Theorem~\ref{thm:finite-correction}.
\end{enumerate}
With uniform constants in these estimates,
\begin{equation}\label{eq:assembly-excess}
 D_{u_0}\le e_1+9\bar e_2+8V_M+O_\varepsilon(y\mathcal B).
\end{equation}
\end{proposition}

\begin{proof}
Average \eqref{eq:loss-second-error} and use $\mu^M_u\ge-e_7$.
Corollary~\ref{cor:bessel-mean} gives
$\langle\mu^N-h\rangle\le\mu^C+e_3$ with $e_3\le8U+12b$. Hence
\[
 \Gamma\le \mu^C+\langle h+E\rangle-\mathcal L+\bar e_2+e_3+e_7 .
\]
Use \eqref{eq:acct-height} for each word. In
$E_u-A\log T_u$ the terms $mAg$ from \eqref{eq:acct-local}
and \eqref{eq:acct-prime} cancel. The identity for $\mathcal M$ in
Proposition~\ref{prop:accounting} and \eqref{eq:loss-first-error} give
\[
 D_{u_0}\le e_1+\bar e_2+e_3+e_4+e_5+Ae_6+e_7 .
\]
By Lemma~\ref{lem:assembly-spread},
$e_3\le8V_M+8\bar e_2+16e_5+12b$.
Sections~\ref{sec:heights} and~\ref{sec:finite-correction} give
$e_4+Ae_6+e_7\ll_\varepsilon y\mathcal B$ and
$e_5\ll_\varepsilon\mathcal B$. Since $b\le\mathcal B$, this proves
\eqref{eq:assembly-excess}.
\end{proof}

\paragraph{An alternative mean-slope input.}
If the other gain, source and height inequalities have been established
for the original family and the other Bessel and cutoff hypotheses
are retained, Proposition~\ref{prop:period-denominator-average} and its
scalar-row variant allow the mean-slope step to use word-dependent
denominators $D_u$.
With the notation of that proposition, translating back to the original
$h_u$ adds at most $9\bar c-8c_{\min}$ to
\eqref{eq:assembly-excess}. Thus $\bar c=o(b)$ is a sufficient
negligible-cost condition. All other accounting terms retain their
original normalization. This alternative does not prove the
denominator estimate or any of the remaining chain, frame or endpoint
comparisons for the prescribed family.

To estimate $V_M$ on the base lattices, write
\[
 V_M\le \max_u\mu_{\max}(M_u)-\langle\mu^M_u\rangle
                 +A\log\frac{\max_u T_u}{\min_u T_u}.
\]
By Proposition~\ref{prop:height-base}, $-\langle\mu^M_u\rangle\le e_7$.
Since $A\le2N_Ta$ and $\log k\ll_\varepsilon y$, it suffices for
$V_M\ll_\varepsilon y\mathcal B$ to verify the bounds of
\citep[Definition~18.9, (K7)--(K8)]{adriazola-grh} on these words,
\[
 \max_u\mu_{\max}(M_u)\ll_\varepsilon N_Ta y,\qquad
 \log\frac{\max_u T_u}{\min_u T_u}\ll_\varepsilon\log k .
\]

\subsection{Selection on the windows}

Suppose the base frame at $u_0$ has rank $D_0\ge\max_j d_j^0$,
$D_0>0$. If its local index comparison holds, Theorem~\ref{thm:twisted-defect}
with Hypothesis~\ref{hyp:twisted-residual} bounds the twisted term
$D_0^{-1}\sum_j H_j$ in Hypothesis~\ref{hyp:height-norms} by
$O_\varepsilon(y\mathcal B)$. The Bessel cutoff still requires a
bound on $s+\mu_{\max}(R)$ in this metric.

\begin{hypothesis}[\openinput]\label{hyp:assembly-windows}
For arbitrarily small $0<\varepsilon<1/52$ there exist
$J\ge3/2$, $K>0$ and $0<\delta_\varepsilon<1$ with the following
property for arbitrarily large $Y$. There are primes $q,d$ of Proposition~\ref{prop:field-window}
on one window $[U_0,U_0^J]\supseteq[Y,Y^{3/2}]$ and a mesh
\[
 Y=x_0<\cdots<x_\nu=Y^{3/2},\qquad
          x_{i+1}-x_i\le Kx_i^{3/4}.
\]
At each $x_i$ there are ranks satisfying \eqref{eq:arithmetic-ranks},
a set $T_2\subset H$ with $|T_2|\ge n x_i^{-\varepsilon}$,
a set $T$ as in Lemma~\ref{lem:weight-count}, and a shift $u_0$.
For these choices there are weights, nodes and lifts as in
\S\ref{sec:theta-chain} for which
Proposition~\ref{prop:assembly-excess} applies and
\begin{equation}\label{eq:assembly-budget}
 e_1+9\bar e_2+8V_M
       \le (1-\delta_\varepsilon)x_i^{5\varepsilon}\mathcal B .
\end{equation}
All constants are uniform over the points and windows for this
$\varepsilon$. Every word and every lift at $x_i$ uses this $T,u_0$.
\end{hypothesis}

Any uniform estimate $e_1+\bar e_2+V_M\ll_\varepsilon
x^{4\varepsilon}\mathcal B$ suffices for \eqref{eq:assembly-budget}
at large $x$. Lemma~\ref{lem:weight-count} gives the translate counts at all shifts.
If the lattices come from a family indexed by shifts, one selects
the index $u_0$ after constructing that family
\citep[\S18.7]{adriazola-grh}. The simultaneous lattice, period
and height estimates on such meshes are open.

\begin{theorem}\label{thm:assembly-rh}
Hypothesis~\ref{hyp:assembly-windows} implies
Hypothesis~\ref{hyp:one-shift-arithmetic} and the Riemann hypothesis.
\end{theorem}

\begin{proof}
Fix $\varepsilon$ from Hypothesis~\ref{hyp:assembly-windows}. Let
$C_\varepsilon$ bound the error in \eqref{eq:assembly-excess} in
units of $x^\varepsilon\mathcal B$. Choose $Y$ large enough that
$C_\varepsilon Y^{-4\varepsilon}\le\delta_\varepsilon$.
At every mesh point, \eqref{eq:assembly-excess}
and \eqref{eq:assembly-budget} give $D_{u_0}\le x^{5\varepsilon}\mathcal B$.
Lemma~\ref{lem:chain-nodes} gives \eqref{eq:weight-ranges}, completing
Hypothesis~\ref{hyp:one-shift-arithmetic}.

With $Q=N_Tmb\delta_2/2$, Proposition~\ref{prop:one-shift} and
the estimate $\mathcal B/Q\ll x^{1/2+3\varepsilon}$ of
\S\ref{sec:arithmetic-bound} give $G(x)\ll_\varepsilon
x^{1/2+8\varepsilon}$ on the mesh by \eqref{eq:field-window-bound}
and \eqref{eq:field-projection}. The interpolation in
Theorem~\ref{thm:arithmetic-rh} extends this bound to $[Y,Y^{3/2}]$.
If $\operatorname{Re}\rho>1/2$ for a nontrivial zero of $\zeta$,
choose $\varepsilon$ with $8\varepsilon<\operatorname{Re}\rho-1/2$.
Theorem~\ref{thm:window-criterion} contradicts the bound on arbitrarily
large windows. The functional equation gives $\operatorname{Re}\rho=1/2$
for every nontrivial zero.
\end{proof}
